\documentclass{article}

% if you need to pass options to natbib, use, e.g.:
%     \PassOptionsToPackage{numbers, compress}{natbib}
% before loading neurips_2024


% ready for submission
% \usepackage{neurips_2024}


% to compile a preprint version, e.g., for submission to arXiv, add add the
% [preprint] option:
%     \usepackage[preprint]{neurips_2024}


% to compile a camera-ready version, add the [final] option, e.g.:
\usepackage[final]{neurips_2024}

\makeatletter
\renewcommand{\@noticestring}{CSE 5100 Spring 2026 Progress Report}
\makeatother

% to avoid loading the natbib package, add option nonatbib:
%    \usepackage[nonatbib]{neurips_2024}


\usepackage[utf8]{inputenc} % allow utf-8 input
\usepackage[T1]{fontenc}    % use 8-bit T1 fonts
\usepackage{hyperref}       % hyperlinks
\usepackage{url}            % simple URL typesetting
\usepackage{booktabs}       % professional-quality tables
\usepackage{amsfonts}       % blackboard math symbols
\usepackage{amsmath}
\usepackage{nicefrac}       % compact symbols for 1/2, etc.
\usepackage{microtype}      % microtypography
\usepackage{xcolor}         % colors


\title{AC-DPO: Adaptive Capacity Direct Preference Optimization}


% The \author macro works with any number of authors. There are two commands
% used to separate the names and addresses of multiple authors: \And and \AND.
%
% Using \And between authors leaves it to LaTeX to determine where to break the
% lines. Using \AND forces a line break at that point. So, if LaTeX puts 3 of 4
% authors names on the first line, and the last on the second line, try using
% \AND instead of \And before the third author name.


\author{
  Aaron Chen \\
  % \thanks{Use footnote for providing further information
  %   about author (webpage, alternative address)---\emph{not} for acknowledging
  %   funding agencies.} \\
  % Department of Computer Science \\
  Washington University in St. Louis \\
  % Pittsburgh, PA 15213 \\
  \texttt{aaronc@wustl.edu} \\
  % examples of more authors
  \And
  Sherlock Zhang \\
  Washington University in St. Louis \\
  \texttt{jiaxi.zhang@wustl.edu} \\
  % \AND
  % Coauthor \\
  % Affiliation \\
  % Address \\
  % \texttt{email} \\
  % \And
  % Coauthor \\
  % Affiliation \\
  % Address \\
  % \texttt{email} \\
  % \And
  % Coauthor \\
  % Affiliation \\
  % Address \\
  % \texttt{email} \\
}

\begin{document}

\maketitle

% \begin{abstract}
%   With the abcde abcde abcde abcde abcde abcde abcde abcde abcde abcde abcde abcde abcde abcde abcde abcde abcde abcde abcde abcde abcde abcde abcde abcde abcde abcde abcde abcde abcde abcde abcde abcde abcde abcde abcde abcde abcde abcde abcde abcde abcde abcde abcde abcde abcde abcde abcde abcde abcde abcde abcde abcde abcde abcde abcde abcde abcde abcde abcde abcde abcde abcde abcde abcde
% \end{abstract}


\vspace{-2.5em}
\section{Brief Recap}
After some more literature review and trial implementation, we decided to stick to our original proposed research topic: Adaptive Capacity Direct Preference Optimization (AC-DPO), which combines curriculum learning over data with dynamic model capacity. Standard DPO is using a fixed LoRA rank throughout the entire training process, which is inefficient to treat low and high difficulty preference pairs the same.

The mismatch can be addressed by the extension we proposed. Our methods divide the training process into two stages: in early stage the model is trained on easier preference pairs with low-rank LoRA adapter; in the later stage, the model is trained with higher model capacity by switching to a high-rank LoRA adapter. The aim of the design is to align model capacity with task difficulty, so that the training efficiency and accuracy can be improved.

\section{Implementation Progress}
By now, we have completed the environment configuration and implemented a working prototype of the AC-DPO pipeline as we planned in the proposal, achieving staged LoRA-rank expansion and adapter switching. We also verified that the model's proper functionality by testing it through a toy problem setting.

\subsection{Training Pipeline}
We firstly implemented the full DPO training pipeline. Referencing the papers we examined in the proposal, we are using the Hugging Face platform. Here is some detailed choice of our implementation:
\begin{itemize}
    \item Base model: GPT-2
    \item Reference Model: frozen copy of the base model
    \item Training framework: \texttt{trl.DPOTrainer}
    \item Parameter efficient fine Tuning: LoRA from \texttt{peft} library
\end{itemize}

The training implementation for now only consists of two separate stages. The sequential stages have their corresponding curriculum level. At each stage, the DPO optimization is performed on that corresponding subset of preference data.

\subsection{Toy Curriculum Construction}
In our original proposal, the curriculum is intended to be constructed by a preference evaluation function, such as reward margins or ranking scores \cite{pattnaik2024currydpo} from a learned reward model. This will skip the process of categorizing the easy and hard preference pairs manually, and is feasible to achieve a large dataset partition.

But for now, this component hasn't been implemented. We instead implemented the Toy Curriculum Dataset with manually defined difficulty levels in order to test our pipeline and make sure it is working. The difficulty level is designed as:
\begin{itemize}
    \item Easy: large difference between preferred and rejected responses
    \begin{itemize}
        \item prompt: Translate to English: Bonjour
        \item chosen: Hello
        \item rejected: I am a banana
    \end{itemize}
    \item Hard: Very subtle difference in reasoning or formatting
    \begin{itemize}
        \item prompt: Write a python function to add two numbers
        \item chosen: \texttt{def add(a, b): \textbackslash n    return a + b}
        \item rejected: \texttt{def sum(a, b): return a+b}
    \end{itemize}
\end{itemize}

The hard-coded approach simplifies the preference pair difficulty, so that we don't have to introduce additional complexity during our validation stage.

\subsection{LoRA Capacity Expansion}
We implemented this unique feature that enables model capacity to increase while the curriculum training proceeds. The two stages have the following capacity details:
\begin{itemize}
    \item Stage 1 (easy)
    \begin{itemize}
        \item LoRA rank: r = 8
        \item Trainable parameters: \textasciitilde 294k
    \end{itemize}
    \item Stage 2 (hard)
    \begin{itemize}
        \item LoRA rank: r = 64
        \item Trainable parameters: \textasciitilde 2.35M
    \end{itemize}
\end{itemize}
The implementation is achieved with the \texttt{peft} adapter mechanism. It allows us to initialize a low-rank adapter at first, add a higher-rank adapter, then during the transition from easy data to hard data training, the active adapter is switched to the corresponding adapter.

\section{Preliminary Results}
We tested the AC-DPO prototype on the toy curriculum dataset and here is the training results summarized in the table.

\begin{table}[h]
\centering
\caption{Detailed Training Metrics Across Stages}
\begin{tabular}{llllll}
\toprule
\textbf{Phase} & \textbf{LoRA Rank} & \textbf{Trainable Params} & \textbf{Loss} & \textbf{Reward Margin} & \textbf{Accuracy} \\
\midrule
Stage 1 (start) & $r = 8$  & \textasciitilde{}294K  & 0.6931 & 0.0000 & 0.0 \\
Stage 1 (end)   & $r = 8$  & \textasciitilde{}294K  & 0.6784 & 0.0297 & 1.0 \\
Stage 2 (start) & $r = 64$ & \textasciitilde{}2.35M & 0.6931 & 0.0000 & 0.0 \\
Stage 2 (early) & $r = 64$ & \textasciitilde{}2.35M & 0.6891 & 0.0081 & 1.0 \\
Stage 2 (end)   & $r = 64$ & \textasciitilde{}2.35M & 0.6816 & 0.0233 & 1.0 \\
\bottomrule
\end{tabular}
\label{tab:training_metrics}
\end{table}

From the results above, the dynamic capacity model is correctly implemented. DPO training pipeline functions correctly in both stages. This confirms that dynamic LoRA expansion is feasible in LLM setting.

In both stages, the training loss decreases consistently, reward margin increases over the training epochs, and the final reward accuracy converges to 1.0 early and remains stable until the training is over. More importantly, at the beginning of Stage 2, the training loss resets to approximately 0.693. This is expected, as Stage 2 introduces both a new dataset and a newly initialized higher-rank adapter. Therefore, the loss values across stages are not directly comparable. After the first few updates in Stage 2, the loss decreases steadily and the reward margin becomes positive, indicating that optimization proceeds normally under the new capacity configuration.

\section{Interpretation}
The preliminary results show feasibility of the proposed AC-DPO method at the toy prototype level, and there are more work to finish the entire model.

First, we conclude that dynamic LoRA expansion can be integrated into multi-stages DPO pipelines. The model runs first in low-rank adapter then switched to a higher-rank adapter on different dataset, and both stages finish with solid convergence. This is a result that we expected, as in previous papers researchers have applied similar techniques in other machine learning areas.

For the optimization behavior, the loss in both stages decreases over training, and the reward margin becomes positive. These results indicate that the model is behaving to assign higher probability to chosen responses under DPO objectives. But this conclusion only works in the stage level, and not across the entire learning process. The two training stages are separate, and the dataset they learn from is separate. Moreover, we didn't implement a baseline model with fixed-capacity and compare the results with AC-DPO, so we can't conclude that it has better performance than the traditional training model yet. We only validate that the pipeline is working correctly here.

At the start of stage 2, the loss increases back to 0.693. However, it is an expected behavior and should not be interpreted as instability during training. The reason behind is that the model is not continuing the same optimization problem. It starts a new problem on a separate dataset. Therefore we only need to focus on the improvement that happens afterward.

\section{Limitation}
The current prototype has huge limitations that we need to work on later.

\subsection{Knowledge Transfer}
Our current implementation has the drawback that the stage two doesn't explicitly inherit or reuse the learning results from Stage 1. It is starting a new optimization process. While this is enough for a pipeline validation, we still need to implement the transition so that we can compare the performance with the baseline fixed-LoRA learning model. We plan to solve this by merging the $r=8$ adapter weights into the base model before initializing the $r=64$ adapter for Stage 2 for continuous knowledge accumulation.

\subsection{Preference Evaluation}
The current toy Curriculum dataset is manually crafted and hard-coded into our code. We still need to implement the preference evaluation function to distinguish whether the given preference pair is easy or hard, or more precisely how hard that pair is quantitatively. This should be feasible as there are plenty of past research providing clear methodology to achieve this goal.

\subsection{Toy-dataset}
Current dataset is too easy. The reward reaches 1.0 in both stages quickly, indicating that the task is still too easy for both stages with different capacities. Future implementation needs larger dataset to test the performance.

\section{Plan for Next Weeks}

\begin{enumerate}
    \item April 6 - April 12 (Weeks 3):
    \begin{enumerate}
        \item Knowledge Transfer (Limitation 5.1): Solve the learning disconnection between learning stages. Before initializing the high-rank adapter of Stage 2, we will implement combined adapter weights, like combining Stage 1 low-rank adapter to the baseline model for consistent learning.
        \item Preference Evaluation (Limitation 5.2 and 5.3): Classify the full-scale preference dataset into "Easy" and "Hard" curriculum types through programmatic calculation of marginalizing rewards. Get rid of the manually handled toy-dataset.
    \end{enumerate}
    
    \item April 13 - April 19 (Weeks 4):
    \begin{enumerate}
        \item Baseline Model: Construct and train a standard, fixed-size (like $r=64$) DPO model on the full-scale dataset.
        \item AC-DPO Pipeline: Run our finalized two-stage AC-DPO pipeline on the full dataset using WashU Academic Compute Cluster.
        \item Presentation Prep: Extract early metrics for visualization like loss curve, marginalized rewards, and checkpoint accuracy for the slides due on April 19.
    \end{enumerate}
    
    \item April 20 - May 3 (Weeks 5-6):
    \begin{enumerate}
        \item Benchmark: Compare the AC-DPO model and the baseline model strictly for metrics like alignment quality, training efficiency, and computation costs.
        \item Documentation: Gather all experiment data and methods to analyze the conclusion for the final report due on May 3. 
    \end{enumerate}
\end{enumerate}

\begin{thebibliography}{4}

\bibitem{croitoru2024curriculum}
Croitoru, F.-A., Hondru, V., Ionescu, R. T., Sebe, N., \& Shah, M. (2024). Curriculum direct preference optimization for diffusion and consistency models. {\it arXiv preprint arXiv:2405.13637}.

\bibitem{deng2026drlora}
Deng, G., Li, B., Chen, R., Wang, H., Wen, L., \& Song, L. (2026). DR-LoRA: Dynamic rank LoRA for mixture-of-experts adaptation. {\it arXiv preprint arXiv:2601.04823}.

\bibitem{pattnaik2024currydpo}
Pattnaik, P., Maheshwary, R., Ogueji, K., Yadav, V., \& Madhusudhan, S. T. (2024). Curry-DPO: Enhancing alignment using curriculum learning \& ranked preferences. {\it arXiv preprint arXiv:2403.07230}.

\bibitem{croitoru2026curriculumdpo}
Croitoru, F.-A., Hondru, V., Ionescu, R. T., Sebe, N., \& Shah, M. (2026). Curriculum-DPO++: Direct preference optimization via data and model curricula for text-to-image generation. {\it arXiv preprint arXiv:2602.13055}.

\end{thebibliography}

\end{document}